\documentclass{article}
% ============================================================
% Bibliography and citations
% ============================================================

\usepackage{natbib}
\usepackage{csquotes}


% ============================================================
% Mathematics
% ============================================================

\usepackage{
    amsmath,
    amssymb,
    amsthm,
    amsfonts,
    mathrsfs,
    dsfont,
    bm,
    textgreek,
    bbm
}


% ============================================================
% Fonts
% ============================================================

\usepackage{newpxtext}
\usepackage{eulervm}


% ============================================================
% Page layout and document structure
% ============================================================

\usepackage[margin=1in, footskip=0.3in]{geometry}
\usepackage{setspace}
\usepackage[bf,small]{titlesec}
\usepackage{etoc}
\usepackage{chngcntr}
\usepackage{changepage}
\usepackage{afterpage}


% ============================================================
% Figures, captions, and graphics
% ============================================================

\usepackage{graphicx}
\usepackage[bf,skip=8pt]{caption}
\usepackage{subcaption}
\usepackage{tikz}

\graphicspath{{Figures/}}


% ============================================================
% Tables
% ============================================================

\usepackage{booktabs}
\usepackage{tabularx}
\usepackage{dcolumn}
\usepackage{makecell}
\usepackage{multirow}
\usepackage{longtable}
\usepackage[para]{threeparttable}
\usepackage{colortbl}
\usepackage{adjustbox}


% ============================================================
% Landscape and rotated content
% ============================================================

\usepackage{lscape}
\usepackage{pdflscape}
\usepackage{rotating}


% ============================================================
% Lists, comments, notes, and text utilities
% ============================================================

\usepackage{enumerate}
\usepackage{comment}
\usepackage{verbatim}
\usepackage{soul}
\usepackage[multiple,stable]{footmisc}
\usepackage[colorinlistoftodos]{todonotes}


% ============================================================
% Code listings
% ============================================================

\usepackage{listings}


% ============================================================
% Colors
% ============================================================

\usepackage{xcolor}

% General-purpose colors
\definecolor{blue}{rgb}{0,0.08,0.5}
\definecolor{red}{rgb}{0.6,0,0}
\definecolor{green}{rgb}{0,0.5,0.05}

% Code-listing colors
\definecolor{fgreen}{HTML}{698B27}
\definecolor{lavender}{HTML}{E6E6FA}
\definecolor{darkpink}{HTML}{C71585}
\definecolor{darkyellow}{HTML}{B8860B}


% ============================================================
% Hyperlinks and PDF metadata
% ============================================================

\usepackage[
    bookmarks,
    pdftitle={},
    pdfauthor={Banks},
    colorlinks=true,
    linkcolor=black,
    citecolor=black,
    urlcolor=black
]{hyperref}


% ============================================================
% Theorem environments
% ============================================================

\newtheorem{prediction}{Prediction}
\newtheorem{assumption}{Assumption}
\newtheorem{proposition}{Proposition}
\newtheorem{corollary}{Corollary}
\newtheorem{definition}{Definition}
\newtheorem{lemma}{Lemma}

% Proposition numbering resets by section
\newtheorem{prop}{Proposition}[section]


% ============================================================
% Mathematical commands
% ============================================================

% Formatting
\newcommand{\no}{\noindent}
\newcommand{\mb}{\mathbf}

% Expectation, probability, and covariance
\newcommand{\E}{\mathbb{E}}
\newcommand{\R}{\mathbb{R}}
\newcommand{\pr}{\mathrm{Pr}}
\newcommand{\cov}{\mathrm{Cov}}

% Support
\newcommand{\supp}[1]{\mathrm{Supp}\left(#1\right)}

% Partial derivatives
\newcommand{\pd}[2]{\frac{\partial #1}{\partial #2}}
\newcommand{\pdd}[2]{\frac{\partial^2 #1}{\partial #2^2}}

% Indicator function
\newcommand{\ind}[1]{\mathbbm{1}\{#1\}}

% Vertical spacing helper for large fractions
\newcommand{\vh}{
    \vphantom{
        \frac{
            \frac{\pd{a}{a}}{\pd{a}{a}}
        }{
            \frac{\pd{a}{a}}{\pd{a}{a}}
        }
    }
}


% ============================================================
% Custom environments
% ============================================================

% Wide equation extending 2 cm into each margin
\newenvironment{widerequation}
{
    \begin{adjustwidth}{-2cm}{-2cm}
    \begin{equation*}
}
{
    \end{equation*}
    \end{adjustwidth}
}

% Solution environment
\newenvironment{solution}
{
    \par\noindent\textbf{\textit{Solution.}}
}
{
    \hfill\rule{1.1ex}{1.1ex}
}


% ============================================================
% Code-listing style
% ============================================================

\lstset{
    basicstyle=\ttfamily\small\color{fgreen},
    keywordstyle=\color{darkpink},
    commentstyle=\color{gray},
    stringstyle=\color{darkyellow},
    showstringspaces=false,
    breaklines=true,
    frame=single,
    backgroundcolor=\color{lavender}
}


% ============================================================
% Spacing and section formatting
% ============================================================

% ECMA-style line spacing
\linespread{1.5}

\titlespacing*{\paragraph}
    {0pt}
    {.25\baselineskip}
    {1em}

\titlespacing*{\section}
    {0pt}
    {.6\baselineskip}
    {.3\baselineskip}

\titlespacing*{\subsection}
    {0pt}
    {.6\baselineskip}
    {.3\baselineskip}

\titlespacing*{\subsubsection}
    {0pt}
    {.6\baselineskip}
    {.3\baselineskip}

% \setlength{\parskip}{5pt}

\setlength{\footnotesep}{.2in}


% ============================================================
% Subsubsection formatting
% ============================================================

\titleformat{\subsubsection}[runin]
    {\normalfont\itshape\bfseries}
    {\thesubsubsection}
    {1em}
    {}
    [.]


% ============================================================
% Date and title-page utilities
% ============================================================

\newcommand{\monthyear}{%
    \ifcase\month\or
        January\or
        February\or
        March\or
        April\or
        May\or
        June\or
        July\or
        August\or
        September\or
        October\or
        November\or
        December
    \fi
    \ \number\year
}


\begin{document}
    \begin{center}
        \scalebox{1.3}{Economics 880 - PS5}\\
        Aaron Banks \& Arthur Rodrigues
    \end{center}


\section*{Question 1}

We solve the household's dynamic programming problem by backward
induction, taking prices $(r,w,b)=(0.05,1.05,0.20)$ as given.
The terminal condition is $V_{N+1}(a,z)=0$, and agents are
restricted to nonnegative asset holdings.

\subsection*{Retirees}

For a retired agent of age $j \geq J^R$, the Bellman equation is
\begin{equation*}
    V_j(a)
    = \max_{a' \geq 0}
    \left\{
        u^R\left((1+r)a+b-a'\right)
        + \beta V_{j+1}(a')
    \right\},
\end{equation*}
where
\begin{equation*}
    u^R(c)=\frac{c^{\gamma(1-\sigma)}}{1-\sigma}.
\end{equation*}

We solve this problem numerically over the asset grid and obtain
the value function for an agent at model-age 50.

\begin{figure}[htbp]
    \centering
    \includegraphics[width=0.75\textwidth]
        {../output/p1_retiree_value.pdf}
    \caption{Value function of a retiree at model-age 50.}
    \label{fig:retiree_value}
\end{figure}

The value function is increasing in assets, since greater wealth
expands the household's feasible consumption set. The underlying
continuous-choice problem also implies a concave value function,
reflecting diminishing marginal utility of consumption. Our numerical
solution is increasing over the entire asset grid and approximately concave.

\subsection*{Workers}

For a working-age agent, the Bellman equation is
\begin{align*}
    V_j(a,z)
    = \max_{a' \geq 0,\;0\leq l\leq 1}
    \Big\{&
        u^W(c,l)
        +\beta \mathbb{E}
        \left[V_{j+1}(a',z')\mid z\right]
    \Big\},
\end{align*}
subject to
\begin{equation*}
    c=(1+r)a+(1-\theta)we(z,j)l-a'.
\end{equation*}

To reduce the dimensionality of the optimization, we derive
optimal labor supply conditional on the savings choice.
The intratemporal first-order condition is
\begin{equation*}
    \frac{1-\gamma}{\gamma}
    \frac{c}{1-l}
    =(1-\theta)we(z,j).
\end{equation*}

Substituting the budget constraint and rearranging yields
\begin{equation*}
    l_j(a,z,a')
    =\gamma-
    \frac{(1-\gamma)\left[(1+r)a-a'\right]}
    {(1-\theta)we(z,j)},
\end{equation*}
with the solution restricted to $[0,1]$.

We then solve for the optimal savings policy $a'_j(a,z)$.
Figure \ref{fig:worker_savings} plots this policy at model-age 20
for both productivity states, $z^H=3.0$ and $z^L=0.5$.

\begin{figure}[htbp]
    \centering
    \includegraphics[width=0.75\textwidth]
        {../output/p1_worker_savings.pdf}
    \caption{Optimal savings policies at model-age 20 by
    productivity state.}
    \label{fig:worker_savings}
\end{figure}

The numerical results indicate that optimal savings are
weakly increasing in current assets for both productivity
states. Moreover, at every point on the asset grid,
\begin{equation*}
    a'_{20}(a,z^H)\geq a'_{20}(a,z^L).
\end{equation*}
Higher-productivity workers earn greater labor income,
holding their assets fixed, and consequently choose to
save more. Both monotonicity properties are satisfied
in our numerical solution.


\section*{Question 2}

We construct the stationary distribution of agents over age,
productivity, and asset holdings using the optimal decision
rules obtained in Question 1.

First, the stationary mass of agents at age $j$ satisfies
\begin{equation*}
    \mu_{j+1}=\frac{\mu_j}{1+n},
    \qquad j=1,\ldots,N-1.
\end{equation*}

Normalizing total population mass to one gives
\begin{equation*}
    \mu_j =
    \frac{(1+n)^{-(j-1)}}
    {\displaystyle\sum_{k=1}^{N}(1+n)^{-(k-1)}}.
\end{equation*}

Newborn agents have zero initial assets, and their
productivity is drawn from the stationary distribution
of the productivity process. Therefore,
\begin{align*}
    F_1(0,z^H) &= 0.2037\mu_1,\\
    F_1(0,z^L) &= 0.7963\mu_1.
\end{align*}

Given the optimal savings policy $a'_j(a,z)$ and
productivity transition matrix $\Pi$, we recursively
compute the distribution at each subsequent age:
\begin{equation*}
    F_{j+1}(a',z')
    =
    \frac{1}{1+n}
    \sum_{a,z}
    F_j(a,z)\Pi(z,z')
    \mathbf{1}\{a'_j(a,z)=a'\}.
\end{equation*}

Here, the indicator function assigns agents to the
asset holdings implied by their optimal savings policy,
while the transition matrix determines the evolution
of productivity.

Our computed distribution satisfies the normalization
conditions
\begin{equation*}
    \sum_{j=1}^{N}\sum_{a,z}F_j(a,z)=1,
\end{equation*}
and
\begin{equation*}
    \sum_{a,z}F_j(a,z)=\mu_j
    \qquad \text{for all }j.
\end{equation*}

Numerically, total population mass is $1.00000000$,
and the maximum discrepancy between the calculated
and theoretical cohort masses is approximately
$8.67\times 10^{-18}$.

The newborn distribution is
\begin{align*}
    F_1(0,z^H)&=0.00430991,\\
    F_1(0,z^L)&=0.01684823.
\end{align*}

Using the stationary distribution, aggregate capital
and effective labor supply are calculated as
\begin{align*}
    K &= \sum_{j=1}^{N}\sum_{a,z}F_j(a,z)a,\\
    L &= \sum_{j=1}^{J^R-1}
         \sum_{a,z}F_j(a,z)e(z,j)l_j(a,z).
\end{align*}

At the fixed prices from Question 1, we obtain
\begin{equation*}
    K=3.622513,\qquad L=0.324883.
\end{equation*}

Figure \ref{fig:assets_by_age} plots average asset
holdings by age, conditional on survival to that age.

\begin{figure}[htbp]
    \centering
    \includegraphics[width=0.75\textwidth]
        {../output/p2_assets_by_age.pdf}
    \caption{Mean asset holdings by age in the
    stationary distribution.}
    \label{fig:assets_by_age}
\end{figure}

Mean asset holdings exhibit a hump-shaped life-cycle
profile. Agents accumulate assets during their
working years to finance retirement consumption,
and subsequently decumulate their wealth after
retirement. This pattern reflects the life-cycle
savings motive in the model.



\section*{Question 3}

We solve for stationary general equilibrium under six specifications,
using the guess-and-verify algorithm. Table \ref{tab:ss_results}
summarizes the results.

\begin{table}[htbp]
    \centering
    \scriptsize
    \setlength{\tabcolsep}{4pt}
    \caption{Steady-state effects of eliminating Social Security}
    \label{tab:ss_results}
    \begin{tabular}{lrrrrrr}
    \hline\hline
     & \multicolumn{2}{c}{Benchmark}
     & \multicolumn{2}{c}{No risk}
     & \multicolumn{2}{c}{Exogenous labor} \\
     & With SS & Without SS
     & With SS & Without SS
     & With SS & Without SS \\
    \hline
    Capital $K$
     & 3.3584 & 4.5984 & 1.0475 & 1.2975 & 7.3553 & 10.4876 \\
    Labor $L$
     & 0.3432 & 0.3651 & 0.1601 & 0.1683 & 0.7542 & 0.7543 \\
    Wage $w$
     & 1.4548 & 1.5931 & 1.2584 & 1.3350 & 1.4530 & 1.6509 \\
    Interest $r$
     & 0.0236 & 0.0112 & 0.0482 & 0.0374 & 0.0238 & 0.0068 \\
    Pension benefit $b$
     & 0.2251 & 0.0000 & 0.0909 & 0.0000 & 0.4941 & 0.0000 \\
    Total welfare $W$
     & -35.7679 & -37.4050 & -45.0110 & -45.1519 & -23.0182 & -25.7607 \\
    $\mathrm{cv}(\mathrm{wealth})$
     & 1.6706 & 1.4946 & 0.7953 & 0.9113 & 1.6565 & 1.4177 \\
    \hline\hline
    \end{tabular}
\end{table}

\subsection*{Benchmark Economy}

The benchmark economy is dynamically efficient, since the equilibrium
interest rate exceeds population growth, $r=0.0236>n=0.011$.
Thus, the implicit return on Social Security is below the return
on private capital.

Eliminating Social Security increases aggregate capital by $36.93\%$
and labor supply by $6.39\%$, raising output by approximately
$16.51\%$. Without pension benefits, households save more to
finance retirement, while eliminating the payroll tax increases
the incentive to work. The resulting increase in the capital-labor
ratio raises wages and lowers the interest rate.

Despite these efficiency gains, aggregate steady-state welfare
falls by $1.6370$. This suggests that the insurance and
redistributive benefits of Social Security outweigh its
distortionary costs under our welfare measure. However, newborn
welfare increases by $1.6394$, with gains for both productivity
types. This highlights the distributional consequences of the
reform: future generations can benefit from higher capital
accumulation, while older generations depend more heavily on
pension benefits. The coefficient of variation of wealth also
falls from $1.6706$ to $1.4946$, indicating lower relative
wealth inequality.

\subsection*{Economy Without Idiosyncratic Risk}

With Social Security, removing idiosyncratic productivity risk
reduces aggregate capital from $3.3584$ to $1.0475$. This is
consistent with a precautionary savings motive: when households
face uncertain labor income, they accumulate assets as a buffer
against adverse productivity shocks. However, the experiment
also lowers average productivity by setting both productivity
states to $0.5$, so the decline in capital cannot be attributed
entirely to the removal of risk.

Eliminating Social Security in this economy increases capital
by $23.87\%$, while aggregate welfare falls by only $0.1410$,
compared with $1.6370$ in the benchmark. This smaller welfare
loss is consistent with Social Security providing valuable
insurance against idiosyncratic income risk, although the
change in mean productivity prevents us from isolating this
mechanism.

Importantly, these welfare comparisons concern stationary
equilibria rather than the transition following a reform.
In particular, they do not account for the losses of existing
generations who have paid Social Security taxes but would
no longer receive their expected benefits. Consequently,
the steady-state welfare differences should not be interpreted
as the welfare effects of an immediate reform.

\subsection*{Exogenous Labor Supply}

Setting $\gamma=1$ eliminates the labor-supply distortion
associated with Social Security. Consistent with this,
aggregate labor is effectively unchanged following the
reform, compared with the $6.39\%$ increase in the benchmark.

Social Security nevertheless continues to crowd out private
capital accumulation. Eliminating the program increases
capital by $42.58\%$, while aggregate welfare falls by
$2.7425$.

Thus, even without the labor-supply distortion, Social
Security reduces capital accumulation but improves aggregate
stationary welfare. The results suggest that the case for
Social Security is stronger when labor supply is exogenous,
as its insurance and redistributive benefits are obtained
without distorting labor supply. Newborn welfare, however,
still increases following the reform.


\end{document}